% Part: normal-modal-logic
% Chapter: axioms-systems
% Section: introduction

\documentclass[../../../include/open-logic-section]{subfiles}

\begin{document}

\olfileid{nml}{prf}{int}

\olsection{مقدمه}

برای زبان پایهٔ موجهات، معناشناسی‌ای برحسب مدل‌های موجهاتی و مفهومی از
معتبر بودن !!a{formula} داریم: صادق‌بودن در همهٔ جهان‌های همهٔ مدل‌ها؛
یا معتبر بودن نسبت به رده‌ای از مدل‌ها یا قاب‌ها: صادق‌بودن در همهٔ
جهان‌های همهٔ مدل‌های رده یا همهٔ مدل‌های مبتنی بر قاب. منطق معمولاً
چنین توصیف‌های معناشناختی اعتبار را به مفهومی نظریه‌اثباتی از
!!{derivability} پیوند می‌دهد. هدف آن است که مفهومی از !!{derivability}
در دستگاهی تعریف کنیم چنان‌که !!a{formula} !!{derivable} باشد اگر و
تنها اگر معتبر باشد.

ساده‌ترین و از نظر تاریخی قدیمی‌ترین دستگاه‌های !!{derivation}،
دستگاه‌های !!{derivation} به‌اصطلاح هیلبرتی یا اصل‌موضوعی‌اند. ساختن دستگاه‌های
!!{derivation} هیلبرتی برای بسیاری از منطق‌های موجهات نسبتاً آسان است:
آن‌ها به‌عنوان موضوع‌های مطالعهٔ فرانظری ساده‌اند (برای نمونه، برای
اثبات صحت و تمامیت). اما اثبات !!{formula}ها در آن‌ها بسیار دشوارتر از،
مثلاً، دستگاه‌های استنتاج طبیعی است.

در دستگاه‌های !!{derivation} هیلبرتی، !!a{derivation} از !!a{formula}
دنباله‌ای از !!{formula}هاست که از اصل‌های موضوع معینی آغاز می‌شود و از
راه چند قاعدهٔ استنتاج به !!{formula} مورد نظر می‌رسد. چون می‌خواهیم
دستگاه !!{derivation} با معناشناسی منطبق باشد، باید تضمین کنیم مجموعهٔ
فرمول‌های !!{derivable} در همهٔ مدل‌ها صادق‌اند (یا در همهٔ مدل‌هایی
که همهٔ اصل‌ها در آن‌ها صادق‌اند). نخست برخی خواص منطق‌های موجهات را
جدا می‌کنیم که برای این کار لازم‌اند: منطق‌های موجهات «عادی». برای
منطق‌های موجهات عادی تنها دو قاعدهٔ استنتاج لازم است: وضع مقدم و
ضرورت‌بخشی. به‌عنوان اصل، همهٔ نمونه‌های جانشینی همان‌گویی‌ها و، بسته
به منطق موجهاتی مورد نظر، تعدادی اصل موجهاتی را می‌گیریم. حتی اگر تنها
به ردهٔ همهٔ مدل‌ها علاقه‌مند باشیم، باید همهٔ نمونه‌های جانشینی
\Ax{K}\iftag{notprvDiamond}{}{ و~$\Ax{Dual}$} را نیز اصل به شمار آوریم.
همین‌ها منطق موجهات عادی کمینه~$\Log K$ را تولید می‌کنند.

\begin{defn}
قاعدهٔ \emph{وضع مقدم} طرحوارهٔ استنتاجی زیر است:
\begin{prooftree}
\AxiomC{$!A$}
\AxiomC{$!A \lif !B$}
\RightLabel{\MP}
\BinaryInfC{$!B$}
\end{prooftree}
می‌گوییم !!a{formula}~$!B$ با وضع مقدم از !!{formula}های~$!A$ و~$!C$
\emph{نتیجه می‌شود} اگر و تنها اگر~$!C \ident !A \lif !B$.
\end{defn}

\begin{defn}
قاعدهٔ \emph{ضرورت‌بخشی} طرحوارهٔ استنتاجی زیر است:
\begin{prooftree}
\AxiomC{$!A$}
\RightLabel{\Nec}
\UnaryInfC{$\Box !A$}
\end{prooftree}
می‌گوییم !!{formula}~$!B$ با ضرورت‌بخشی از !!{formula}~$!A$ نتیجه
می‌شود اگر و تنها اگر~$!B \ident \Box !A$.
\end{defn}

\begin{defn}
یک \emph{!!{derivation}} از مجموعهٔ اصل‌های~$\Sigma$ دنباله‌ای از
!!{formula}های~$!B_1$، $!B_2$، \dots، $!B_n$ است که هر~$!B_i$ یکی از
موارد زیر باشد:
\begin{enumerate}
\item یک نمونهٔ جانشینی از همان‌گویی؛ یا
\item یک نمونهٔ جانشینی از !!a{formula} در~$\Sigma$؛ یا
\item با وضع مقدم از دو !!{formula}~$!B_j$ و~$!B_k$، که~$j,k<i$،
  نتیجه شود؛ یا
\item با ضرورت‌بخشی از !!a{formula}~$!B_j$، که~$j<i$، نتیجه شود.
\end{enumerate}
اگر چنین !!a{derivation}ای با~$!B_n \ident !A$ وجود داشته باشد،
می‌گوییم~$!A$ \emph{از~$\Sigma$ !!{derivable} است}، به نماد
$\Sigma \Proves !A$.
\end{defn}

با این تعریف روشن خواهد شد که مجموعهٔ !!{derivable}ها یک منطق
موجهات عادی تشکیل می‌دهد و هر !!{derivable} !!{formula} در هر مدلی که
همهٔ اصل‌ها در آن صادق‌اند، صادق است. این خاصیت !!{derivation}ها
\emph{صحت} نامیده می‌شود. اثبات عکس آن، یعنی \emph{تمامیت}، دشوارتر است.

\end{document}
